\documentclass[12pt]{article}
\usepackage[utf8]{inputenc}
\usepackage[T1]{fontenc}
\usepackage{hyperref}
\title{Wangaaybuwan-Ngiyambaa Language Resources for K-12 Curriculum Development}
\author{Compiled by Claude Sonnet 5.5 (claude-sonnet-5-5) for Andr\'as Kornai}
\date{}
\begin{document}
\maketitle

\section{Introduction}
Wangaaybuwan-Ngiyambaa (ISO 639-3 wyb; Glottocode wang1291) is a Pama-Nyungan language of New South Wales, Australia. It is also written Wangaaypuwan Ngiyampaa. Sources disagree on its status. Glottolog marks it as extinct, while Wikipedia gives an estimate of 11--50 speakers for Ngiyampaa generally. Reference materials are mostly academic and published by the community and its collaborators.

\section{Primary Educational Resources}
I found no school textbook or primer in the search results. The main community-oriented resources are the dictionary and word book below.
\begin{itemize}
\item Donaldson, \emph{Ngiyampaa Wordworld: Thipingku Yuwi, Maka Ngiya, Names of Birds and Other Words} (Aboriginal Studies Press, 2020), which includes an English-to-Ngiyampaa finderlist: \url{https://shop.aiatsis.gov.au/products/ngiyampaa-wordworld}
\end{itemize}

\section{Linguistic and Reference Materials}
\begin{itemize}
\item Donaldson, \emph{Ngiyambaa: the Language of the Wangaaybuwan} (1980), the grammar (AIATSIS publication record): \url{https://aiatsis.gov.au/publication/35862}
\item Ngiyampaa language (Wikipedia): \url{https://en.wikipedia.org/wiki/Ngiyampaa_language}
\item Glottolog entry for Wangaaybuwan-Ngiyambaa: \url{https://glottolog.org/resource/languoid/id/wang1291}
\item Endangered Languages Project records: \url{https://endangeredlanguages.com/node/92257} and \url{https://endangeredlanguages.com/node/38813}
\end{itemize}

\section{Religious and Community Resources}
\begin{itemize}
\item Ngiyampaa Wangaaypuwan community website, which hosts a searchable dictionary and includes map and community context: \url{https://wangaaypuwan.org.au}
\end{itemize}
I saw no scripture resources in the search results.

\section{Digital and Audio Resources}
The community website above provides a searchable online dictionary by English word. I saw no audio or video collections in the search results.

\section{Curriculum Development Context}
The language is severely endangered or dormant, and no school programme was found in the results. A heritage curriculum would have to be built from Donaldson's grammar and Wordworld and from the community's online dictionary, with the Wangaaypuwan community directing the work and its permissions. Pronunciation and audio support would need community involvement.

\end{document}